% Generated from ../chapters/21-relaxation-as-active-release.md
\chapter{Chapter 21 --- Relaxation as Active Release}\label{chapter-21-relaxation-as-active-release}

Relaxation is often misunderstood as doing nothing. In this framework, relaxation is the active removal of constraints that no longer serve the organism.

A guarded joint has fewer available trajectories. A person anticipating pain contracts before movement. Chronic stress narrows attention and breathing. These constraints may once have been protective. When they remain after the original threat, they become a local attractor.

The goal is not maximal looseness. Stability is necessary. The goal is to release surplus control while preserving safety.

A practical session can be organized around five principles:

\begin{enumerate}
\def\labelenumi{\arabic{enumi}.}
\tightlist
\item
  \textbf{Low force.} The movement should remain reversible.
\item
  \textbf{Slow time.} Speed should allow sensation to guide adjustment.
\item
  \textbf{Broad attention.} Notice relations across the body rather than fixating on one point.
\item
  \textbf{No dramatic target.} Do not seek cracks, pain, or extreme range.
\item
  \textbf{After-effect.} Judge the session by durable function afterwards.
\end{enumerate}

The nervous system may initially interpret freedom as danger. The old pattern can return. Repetition provides evidence that another configuration is safe.

Relaxation also applies psychologically. A rigid identity predicts which actions are possible. Allowing uncertainty can release cognitive degrees of freedom. This is why contemplative practice and ambiguous intellectual work can become parts of the same framework.
